\documentclass[11pt,a4paper]{article}
\usepackage{amsmath,amssymb,geometry}
\geometry{margin=2cm}
\allowdisplaybreaks
\title{集成交易与利润最大化}
\author{}
\date{}
\usepackage[UTF8]{ctex}
\begin{document}
\maketitle

保留的配方是单一情景综合模型（1–24），其中由源定义的类比燃料消耗函数被精确替换。所有与情景相关的数据都指向同一个固定情景。

\section*{集合与索引}
\begin{itemize}
\item $T\in\mathbb Z_{>0}$表示周期的数量，而$\mathcal T=\{1,\ldots,T\}$则代表其他数值。请将这些数学符号按原顺序原封不动地复制下来。所有数字都保持不变。
\item $J_t\in\mathbb Z_{>0}$表示周期$t$和$\mathcal J_t=\{1,\ldots,J_t\}$中分段的个数。
\item $\mathcal A=\{(\tau,t)\in\mathcal T^2:\tau\le t\}$表示合同期限和交付期限。
\item $\varepsilon=\{\mathrm{on},\mathrm{off}\}$表示高峰时段和非高峰时段的电力使用情况。提供的映射关系$p:\bigcup_{t\in\mathcal T}\mathcal J_t\to\varepsilon$将高峰状态分配给段索引$j$。
\item $I$是有限生成器集合；$\mathrm{Gas},\mathrm{Coal},\mathrm{Nuc}\subseteq I$则是所提供的、与燃料相关的生成器集合。
\item $\mathcal D$是有限的天数集合。对于每个$d\in\mathcal D$，$t(d)\in\mathcal T$表示其周期，而$j(d)\subseteq\mathcal J_{t(d)}$则包含其两个分段索引。
\end{itemize}
所有求和和约束的索引范围都是整数范围。空范围不会产生任何约束，而对空集合的求和结果为零。

\section*{固定参数和边界数据}
该实例提供了以下参数。除了索引计数和明确指定的离散数据外，系数都是实数标量，其索引方式也如所述。
\begin{description}
\item[$H_e$:] 段落的持续时间，以小时为单位，其中包含$H_{\mathrm{on}}=16$和$H_{\mathrm{off}}=8$。
\item[$P$:] 规范后的电力销售价格。
\item[$\alpha$:] 实数值的流动性限制系数。
\item[$D_{tj},PS_{tj}$:] 分别为$t\in\mathcal T$和$j\in\mathcal J_t$的电力需求与现货电价。
\item[$PP_{\tau te},PG_{\tau t}$:] 分别为$(\tau,t)\in\mathcal A$的电力和天然气远期价格，其中$e\in\varepsilon$表示电力价格。
\item[$PT_t$:] 第$t\in\mathcal T$期所需的利润门槛。
\item[$ML$:] 每日允许的最大损失额。
\item[$CP_t,NP_t$:] $t\in\mathcal T$地区的煤炭和核燃料价格。
\item$q_i,Q_i$: 对于$i\in I$来说，其最低和最高运行发电水平是多少。
\item[$L_i,l_i$:] 以段落单位为单位，为$i\in I$提供了整数形式的最小运行时间以及停机时间长度。
\item[$a_i,b_i$:] 每个$i\in\mathrm{Gas}\cup\mathrm{Coal}\cup\mathrm{Nuc}$对应的仿射型燃料消耗截距和斜率。该截距与承诺状态无关。
\item[$W_{itj},\varpi_{itj}\in\{0,1\}$:] $i\in I$、$t\in\mathcal T$、$j\in\mathcal J_t$的预定和强制可用性指示符；零表示不可用。
\item[$YP_{0te}^{+},YP_{0te}^{-}\ge0$:] $t\in\mathcal T$和$e\in\varepsilon$的初始电力持有量分别等于提供的初始持有量$YL_{0te}$和$YS_{0te}$。
\item[$YG_{0t}\ge0$:] $t\in\mathcal T$的初始气体储量。
\item[$U_{it0}\in\{0,1\}$:] 为每个$i\in I$、$t\in\mathcal T$提供了初始操作状态。这些数值是每个周期固定的边界数据。
\item[$\overline{FP}_{\tau te}^{+},\overline{FP}_{\tau te}^{-},\overline{YP}_{\tau te}^{+},\overline{YP}_{\tau te}^{-}$:] 为$(\tau,t)\in\mathcal A$和$e\in\varepsilon$提供了电力交易和持有成分的非负上限。 
\end{description}

\section*{决策变量与域}
对于$(\tau,t)\in\mathcal A$、$e\in\varepsilon$、$FP_{\tau te}$，它们分别表示有符号的电力远期交易；而$YP_{\tau te}$则表示有符号的电力远期持有量。正数和负数分别表示多头和空头头寸。它们的多头和空头幅度满足以下条件：
\begin{align*}
FP_{\tau te},YP_{\tau te}&\in\mathbb R,
&& (\tau,t)\in\mathcal A,\ e\in\varepsilon,\\
0\le FP_{\tau te}^{+}&\le\overline{FP}_{\tau te}^{+},
&& (\tau,t)\in\mathcal A,\ e\in\varepsilon,\\
0\le FP_{\tau te}^{-}&\le\overline{FP}_{\tau te}^{-},
&& (\tau,t)\in\mathcal A,\ e\in\varepsilon,\\
0\le YP_{\tau te}^{+}&\le\overline{YP}_{\tau te}^{+},
&& (\tau,t)\in\mathcal A,\ e\in\varepsilon,\\
0\le YP_{\tau te}^{-}&\le\overline{YP}_{\tau te}^{-},
&& (\tau,t)\in\mathcal A,\ e\in\varepsilon.
\end{align*}
气体远期购买$FG_{\tau t}$和持有量$YG_{\tau t}$都有各自的领域。
\[
FG_{\tau t},YG_{\tau t}\in\mathbb R_{\ge0}
\qquad ((\tau,t)\in\mathcal A).
\]
剩余的变量包括：
\begin{itemize}
\item $SP_{tj}$: 有符号的现货电力交易，正数表示购买，负数表示销售；
\item $ZP_{te},ZG_t$：总的电力和天然气输送成本；
\item $G_{itj},TG_{tj}$: 个体生成与总生成量；
\item $SG_{tj},SC_{tj},SN_{tj}$: 天然气、煤炭和核燃料的消耗量；
\item $C_{tj}$: 生成成本；
\item $U_{itj}$：运行指示灯，当发电机$i$开启时，该指示灯的值为1。
\end{itemize}
他们的领域是：
\begin{align*}
ZP_{te}&\in\mathbb R
&&t\in\mathcal T,\ e\in\varepsilon,\\
ZG_t&\in\mathbb R
&&t\in\mathcal T,\\
SP_{tj},TG_{tj},SG_{tj},SC_{tj},SN_{tj},C_{tj}&\in\mathbb R
&&t\in\mathcal T,\ j\in\mathcal J_t,\\
G_{itj}&\in\mathbb R
&&i\in I,\ t\in\mathcal T,\ j\in\mathcal J_t,\\
U_{itj}&\in\{0,1\}
&&i\in I,\ t\in\mathcal T,\ j\in\mathcal J_t.
\end{align*}
所有连续变量都是实数值，并且满足以下约束条件。

\section*{目标}
对声明的决策变量进行最大化处理：
\begin{equation}
\max\quad
\sum_{t\in\mathcal T}\sum_{j\in\mathcal J_t}
\left[(D_{tj}P-SP_{tj}PS_{tj})H_{p(j)}-C_{tj}\right]
-\sum_{t\in\mathcal T}\sum_{e\in\varepsilon}ZP_{te}.
\tag{24}
\end{equation}

\section*{交易限制}
\begin{align}
FP_{\tau te}&=FP_{\tau te}^{+}-FP_{\tau te}^{-}
&& (\tau,t)\in\mathcal A,\ e\in\varepsilon,\tag{1}\\
YP_{\tau te}&=YP_{\tau te}^{+}-YP_{\tau te}^{-}
&& (\tau,t)\in\mathcal A,\ e\in\varepsilon,\tag{2}\\
YP_{\tau te}^{+}&=YP_{(\tau-1)te}^{+}+FP_{\tau te}^{+}
&& (\tau,t)\in\mathcal A,\ e\in\varepsilon,\tag{3}\\
YP_{\tau te}^{-}&=YP_{(\tau-1)te}^{-}+FP_{\tau te}^{-}
&& (\tau,t)\in\mathcal A,\ e\in\varepsilon,\tag{4}\\
YG_{\tau t}&=YG_{(\tau-1)t}+FG_{\tau t}
&& (\tau,t)\in\mathcal A,\ e\in\varepsilon,\tag{5}\\
\sum_{t=\tau}^{T}FP_{\tau te}^{+}
&\le\alpha\sum_{t=\tau}^{T}YP_{(\tau-1)te}^{+}
&& \tau\in\mathcal T,\ e\in\varepsilon,\tag{6}\\
\sum_{t=\tau}^{T}FP_{\tau te}^{-}
&\le\alpha\sum_{t=\tau}^{T}YP_{(\tau-1)te}^{-}
&& \tau\in\mathcal T,\ e\in\varepsilon,\tag{7}\\
ZP_{te}&=\frac{J_t}{2}\sum_{\tau=1}^{t}PP_{\tau te}FP_{\tau te}H_e
&&t\in\mathcal T,\ e\in\varepsilon,\tag{8}\\
ZG_t&=\sum_{\tau=1}^{t}PG_{\tau t}FG_{\tau t}
&&t\in\mathcal T.\tag{9}
\end{align}
\begin{equation}
\sum_{j\in\mathcal J_t}
\left[(D_{tj}P-SP_{tj}PS_{tj})H_{p(j)}-C_{tj}\right]
-\sum_{e\in\varepsilon}ZP_{te}\ge PT_t
\qquad t\in\mathcal T.
\tag{10}
\end{equation}

\section*{能量平衡与发电核算}
\begin{align}
YP_{tt,p(j)}+SP_{tj}+TG_{tj}&=D_{tj}
&&t\in\mathcal T,\ j\in\mathcal J_t,\tag{11}\\
YG_{tt}&=\sum_{j\in\mathcal J_t}SG_{tj}
&&t\in\mathcal T,\tag{12}\\
SG_{tj}&=\sum_{i\in\mathrm{Gas}}a_i
       +\sum_{i\in\mathrm{Gas}}b_iG_{itj}
&&t\in\mathcal T,\ j\in\mathcal J_t,\tag{13}\\
SC_{tj}&=\sum_{i\in\mathrm{Coal}}a_i
       +\sum_{i\in\mathrm{Coal}}b_iG_{itj}
&&t\in\mathcal T,\ j\in\mathcal J_t,\tag{14}\\
SN_{tj}&=\sum_{i\in\mathrm{Nuc}}a_i
       +\sum_{i\in\mathrm{Nuc}}b_iG_{itj}
&&t\in\mathcal T,\ j\in\mathcal J_t,\tag{15}\\
C_{tj}&=\frac{ZG_t}{J_t}+NP_tSN_{tj}+CP_tSC_{tj}
&&t\in\mathcal T,\ j\in\mathcal J_t,\tag{16}\\
TG_{tj}&=\sum_{i\in I}G_{itj}
&&t\in\mathcal T,\ j\in\mathcal J_t.\tag{17}
\end{align}

\section*{操作与可用性限制}
\begin{equation}
q_iU_{itj}\le G_{itj}\le Q_iU_{itj}
\qquad i\in I,\ t\in\mathcal T,\ j\in\mathcal J_t.
\tag{18}
\end{equation}
\begin{equation}
\begin{gathered}
U_{itj}-U_{it,j-1}\le U_{it\tau},\\
i\in I,\quad t\in\mathcal T,\quad j\in\mathcal J_t,\quad
\tau\in\{j+1,\ldots,\min(j+L_i-1,J_t)\}.
\end{gathered}
\tag{19}
\end{equation}
\begin{equation}
\begin{gathered}
U_{it,j-1}-U_{itj}\le 1-U_{it\tau},\\
i\in I,\quad t\in\mathcal T,\quad j\in\mathcal J_t,\quad
\tau\in\{j+1,\ldots,\min(j+l_i-1,J_t)\}.
\end{gathered}
\tag{20}
\end{equation}
\begin{align}
U_{itj}&\le W_{itj}
&&i\in I,\ t\in\mathcal T,\ j\in\mathcal J_t,\tag{21}\\
U_{itj}&\le\varpi_{itj}
&&i\in I,\ t\in\mathcal T,\ j\in\mathcal J_t.\tag{22}
\end{align}

\section*{每日亏损限额}
\begin{equation}
\begin{gathered}
\sum_{j\in j(d)}
\left[(D_{tj}P-SP_{tj}PS_{tj})H_{p(j)}-C_{tj}\right]
-\frac{2}{J_t}\sum_{e\in\varepsilon}ZP_{te}+ML\ge0,\\
d\in\mathcal D,\qquad t=t(d).
\end{gathered}
\tag{23}
\end{equation}
\end{document}
